\fichatitulo{45 · CONTEXTO}{ISKO Brasil · 2025}{IA na organização do conhecimento parlamentar}
{\small Viola; Felipe; Sales. \textit{Intelig{\^e}ncia Artificial na Organiza{\c c}{\~a}o do Conhecimento do Parlamento brasileiro: um estudo aplicado aos dados e {\`a}s informa{\c c}{\~o}es legislativas}. ISKO Brasil · 2025. \href{https://isko.org.br/ojs/index.php/iskobrasil/article/download/52/54/685}{Fonte consultada}.\par}

\campo{Problema e objetivo}
Como a Câmara utiliza IA na organização do conhecimento?

\campo{Principal contribuição · síntese interpretativa}
Mapeia iniciativas e resultados esperados a partir de documentos institucionais.

\campo{Método / natureza da evidência}
Pesquisa qualitativa, exploratória, bibliográfica e documental; consulta portais da Câmara e da IPU.

\campo{Resultado ou argumento principal}
Agrupa seis ferramentas em quatro categorias, incluindo transcrição, classificação e recuperação.

\campo{Excertos-chave · seleção literal}
\begin{itemize}
\excerto{1}{abordagem qualitativa com objetivo exploratório}{Seção 2.}
\excerto{2}{para a identificação do objetivo e dos resultados esperados}{Seção 4.}
\end{itemize}

\campo{Limitações e análise crítica}
Análise da ficha: descrição documental e resultados esperados não demonstram impacto causal nem desempenho comparativo de RAG.

\campo{Aproveitamento na dissertação · proposta}
Contextualização: identificar tarefas legislativas e distinguir benefícios esperados de resultados medidos.

\registro{Fonte consultada nesta preparação: Extração web do PDF oficial: resumo e seções 2–4, até recuperação de documentos. Sem PDF local; leitura parcial. Chave: \texttt{martelloteViola2025organizacaoConhecimento}.}
